% ============================================================================
%  Task_W02.tex -- Backtracking, Week 2
%  Uses coa-accessible.sty for the accessible baseline (headings, captions,
%  color palette, page numbers, PDF metadata). Compile twice with pdflatex
%  so "Page X of Y" resolves.
% ============================================================================
\documentclass[11pt]{article}
\usepackage{coa-accessible}

% Plain, language-neutral pseudocode listing style, defined locally so the
% shared style file stays untouched. No keyword highlighting, since this
% pseudocode is not tied to Java or C or C++.
\lstdefinestyle{pseudo}{
  basicstyle=\ttfamily\small\color{coaBodyText},
  frame=single,
  rulecolor=\color{coaNavy},
  breaklines=true,
  showstringspaces=false,
  columns=fullflexible,
  tabsize=2,
  numbers=none
}
\lstset{style=pseudo}

\AssignmentMeta{Task W02: Backtracking}{Week 2 (Choose, Explore, Backtrack)}

\begin{document}

\AssignmentTitleBlock

\begin{center}
\rule{0.95\linewidth}{0.6pt}
\end{center}
\subsection*{A note on how this assignment was developed}
\noindent
This assignment was drafted with AI help (mostly for brainstorming new
problems to solve and formatting with LaTeX). It should be accurate, but
something might still be wrong.
Because this is a new assignment, you are also helping me debug it, and I
will take that into consideration when I review your work. That said,
helping me find rough edges is not a substitute for doing the assignment;
each problem below is solvable, but you may need to make assumptions in
your attempt. If you find an error or
something unclear, please tell me; that helps later versions, though
finding errors is not the point of this assignment.
\begin{center}
\rule{0.95\linewidth}{0.6pt}
\end{center}
\vspace{0.6em}

\section{Important Guidelines}
\begin{itemize}[leftmargin=1.4em]
  \item Exam questions and problems may extend beyond tasks and lab assignments.
  \item No Generative AI: You must complete this assignment using only the
        resources provided in the course (book, lectures, and materials).
        Credit will not be granted for using external features or techniques
        not covered in this course. (See the one exception in the
        Eight-Queens section.)
  \item Ask Me Questions: If you need clarification, please ask me directly.
        Other instructors or advisors may not be aware of the specific
        details and expectations of this assignment.
  \item The assignment instructions do not detail every step; therefore, you
        must read, interpret, and possibly make assumptions. If something is
        unclear, write down the assumption you make, proceed with it, and
        complete the problem. Do not leave a problem unfinished because a
        detail was ambiguous. Not all assumptions are automatically correct,
        so state yours clearly enough that I can follow your reasoning.
\end{itemize}

\section{How to Submit Your Work}
Everything for this assignment goes on your answer sheet: the boards for
Problems 1 and 2, your backtrack notes, your rewritten pseudocode, and your
written answers for Problem 3. There is nothing to draw or submit
separately.

\subsection*{If you find something that needs to be debugged}
If, while working through this assignment, you discover something in the
problem statements that needs to be debugged, describe what you found in
the "Bugs or Issues You Found" section of the answer sheet.

\subsection*{Where to submit}
Commit your full submission, including your answer sheet's \texttt{.tex}
source, its compiled PDF, and your AI conversation if you used the
formatting exception in the Eight-Queens section, to a folder named
\texttt{taskw02} in your git repository.

\section{How These Problems Work}

A backtracking search follows the same three-part pattern every time: choose
an option, explore what happens if you take it, and if that exploration
cannot succeed, undo the choice and backtrack to try the next option. This is
the same calls-and-returns idea from the recursion basics assignment, except
a call here can have several options to try instead of one, and some of
those options fail before reaching a base case.

For Problems 1 and 2, you will not diagram the entire search. Instead you
will fill in the board on your answer sheet with the algorithm's final
result: queen positions for Eight-Queens. Where a problem asks for
backtrack notes, record them in the Backtrack Notes section of your answer
sheet; the point is to show you actually worked through the algorithm
rather than guessed at an answer, without requiring a full diagram of every
call. Problem 3 takes a different approach, using a video instead of hand
execution; see Section 5 for details.

\section{Eight-Queens}

Eight-Queens asks whether $N$ queens can be placed on an $N \times N$ board
so that no two queens share a row, a column, or a diagonal. The book's
version uses an eight by eight board, which is far too large to work through
reliably by hand. Both problems below use a four by four board instead,
small enough to work through completely.

The pseudocode below places one queen per column, choosing which row to try
it in, and stops as soon as one full solution is found.

\begin{lstlisting}
function SolveQueensByColumn(column):
    if column == N:
        return true              // every column has a queen; solved
    for row = 0 to N - 1:
        if placing a queen at (row, column) is safe, given the
                queens already placed in earlier columns:
            place a queen at (row, column)
            if SolveQueensByColumn(column + 1):
                return true      // solution found; stop searching
            remove the queen at (row, column)     // backtrack
    return false                 // dead end: no row worked here
\end{lstlisting}

\subsection{Problem 1: A Smaller Board}
Using the pseudocode above exactly as written, work through
\texttt{SolveQueensByColumn} on a four by four board, starting from column
0. Fill in the board on your answer sheet, labeled with columns 1 through 4
left to right and rows 1 through 4 top to bottom, and place a queen in each
square where the algorithm's first full solution ends up.

For any column where the row you tried first did not work, add one short
note in the Backtrack Notes section of your answer sheet: which column,
which row you tried first, and why it failed.

\subsection{Problem 2: Row by Row}
Before you work through anything, rewrite the pseudocode above so that it
places one queen per row instead of one queen per column, choosing which
column to try the queen in for each row. Write your rewritten pseudocode in
the space provided on your answer sheet for this problem.

Then work through your rewritten version on the same four by four board,
starting from row 0, and fill in the board on your answer sheet the same
way you did for Problem 1, including one short note in the Backtrack Notes
section for any row where your first-tried column did not work.

Once you have solved this problem, and only after you have solved it, you
may use an AI tool solely to format your finished pseudocode. If you do
this, you must explicitly tell the AI not to change or fix your algorithm,
only to format it, and you must submit the full AI conversation alongside
your pseudocode. This exception does not extend to using AI to solve the
problem, or to fix a mistake in your algorithm; that is still covered by
the No Generative AI guideline above.

\section{Knight's Tour}
\label{sec:knights-tour}

A knight's tour asks whether a knight can visit every square on a board
exactly once, moving only in the knight's usual L-shape, never revisiting a
square. Unlike Eight-Queens, a single call in this search has up to eight
options to try, since a knight can move in up to eight directions from any
square, fewer near an edge or a corner.

The pseudocode below tries a knight's eight possible moves in a fixed order
every time, and stops as soon as the entire board has been visited.

\begin{lstlisting}
function SolveKnightsTour(row, col, squaresVisitedSoFar):
    mark (row, col) as visited
    if squaresVisitedSoFar == totalSquares:
        return true                     // every square visited; solved

    rowOffset = [-2, -2, -1, -1,  1,  1,  2,  2]
    colOffset = [-1,  1, -2,  2, -2,  2, -1,  1]

    for i = 0 to 7:
        nextRow = row + rowOffset[i]
        nextCol = col + colOffset[i]
        if 0 <= nextRow < numRows and 0 <= nextCol < numCols
                and (nextRow, nextCol) is not visited:
            if SolveKnightsTour(nextRow, nextCol, squaresVisitedSoFar + 1):
                return true             // solution found; stop searching

    mark (row, col) as unvisited         // backtrack: undo this square so
    return false                         // the caller can try its next move
\end{lstlisting}

\subsection{Worked Example}
Here is what \texttt{SolveKnightsTour} finds on a four row by three column
board (twelve squares), starting from square B1 (column B, row 1). Columns
run A, B, C left to right; rows run 1 through 4 top to bottom. The number in
each square below is the order the knight visited it.

\begin{accessibletable}{Knight's Tour worked example: visit order starting from B1}{tab:knight-demo}{cccc}
\toprule
 & A & B & C \\
\midrule
1 & 10 & 1 & 8 \\
2 & 7 & 4 & 11 \\
3 & 2 & 9 & 6 \\
4 & 5 & 12 & 3 \\
\bottomrule
\end{accessibletable}

Reaching this board was not a straight line. It took 31 total recursive
calls and 19 backtracks to get here. For example, after ten confident moves
in a row, square C4 turned out to be a dead end: every neighboring square
was already visited. The search had to backtrack through four earlier
squares, C4, B2, A4, and C3, before a new path forward finally opened up at
A2's next candidate move. This is exactly the choose, explore, backtrack
pattern from Section 3; it is just not diagrammed here move by move.

Notice that this small twelve-square board already needed 31 calls, nearly
twice the number of squares on the board. Keep that in mind for the
reflection questions.

\subsection{Problem 3: Watch and Answer}
Watch this video on the Knight's Tour puzzle:

\begin{center}
\textit{Knight's Tour} --- Numberphile\\
\url{https://www.youtube.com/watch?v=ab_dY3dZFHM}
\end{center}

\noindent As you watch, notice how many possible paths the knight has to
consider, and how much of solving this by hand really amounts to organized
trial and error, the same choose, explore, backtrack pattern from Section 3.

Then answer the following in a few sentences each on your answer sheet.

\begin{enumerate}[leftmargin=1.4em]
  \item In your own words, explain what backtracking means for a knight's
        tour search. Use the words "dead end" and "undo" somewhere in your
        answer.
  \item Based on what you saw with Eight-Queens, would you expect a
        backtracking knight's tour search starting from a corner square to
        need more backtracking, less, or is it impossible to say without
        trying, compared to starting from a square near the middle of the
        board? Explain your reasoning, thinking about how many legal moves
        a knight has from each type of square.
  \item Suppose the fixed order in which the eight knight moves are tried,
        given in the pseudocode above, was changed to a different fixed
        order, similar to how Problem 2 changed Eight-Queens from column by
        column to row by row. Would you expect a valid tour to still turn
        up just as quickly, or could the order matter here too? Explain
        your reasoning.
  \item A three by three board has a center square that a knight can never
        move to or from anywhere else on that board. Based on this, could a
        backtracking search ever find a complete tour on a three by three
        board, no matter where it starts or what order it tries moves in?
        Explain your reasoning.
\end{enumerate}

\section{Reflection Questions}
Answer these in a few sentences each, after you have completed your work
above.

\begin{enumerate}[leftmargin=1.4em]
  \item Compare your Problem 1 board to your Problem 2 board, along with the
        backtrack notes you wrote for each. Are the boards the same,
        mirrored, or something else entirely? Did one problem need more
        backtrack notes than the other? Explain why, based on the shape of
        the board and how the safety check works.
  \item The Knight's Tour worked example took 31 total calls and 19
        backtracks on a twelve-square board. Your Problem 1 board has
        sixteen squares, more than the knight's board, yet almost certainly
        needed far fewer backtrack notes. What is different about the two
        problems that could explain this?
  \item Look back at your answer to Problem 3, question 2, about corner
        versus central starting squares. Does that reasoning match what you
        actually saw happen across your Eight-Queens work, where changing
        where the search starts (column 0 for Problem 1, row 0 for Problem
        2) still ended up exploring a similar amount? Explain any
        similarities or differences you notice.
  \item Problem 2 asked you to change Eight-Queens from column by column to
        row by row. Suppose you instead kept the queens problem exactly as
        written for Problem 1, and only changed the order that rows are
        tried within each column, for example trying the bottom row first
        instead of the top row. Based on what you saw across Problems 1
        through 3, would you expect that change to behave more like Problem
        2, or more like the difference you saw between Problem 1 and the
        Knight's Tour example? Explain your reasoning.
\end{enumerate}

\end{document}
